\section{The annealing algorithm and its implementation}
\label{sec:algorithm}

We now combine the preceding estimates into a single algorithm.  The
transversal distribution at parameter \(j\) is denoted by
\(\mu_j=f_j/C(j)\).  It does not depend on the defect multipliers.
The underlying Metropolis kernel with parameter \(j\) and multiplier
collection \(w\) is denoted by \(P_{j,w}\).  We use the warm trace
restarts of Section~\ref{sec:transport} to estimate successive partition
function ratios while learning the defect multipliers. Balancing
defect classes by estimated weights during annealing is an important
feature of the permanent approximation algorithm of Jerrum, Sinclair,
and Vigoda~\cite[Section~3]{JSV2004}. Here the observable-variance and trace bounds
needed for the two-lift weights were proved in
Section~\ref{sec:transport}.

The feasibility test and tightening from Section~\ref{sec:scales} are
performed first.  An infeasible instance returns zero.  Throughout this
section the instance is feasible, so \(Z\geq 1\).  Recall that
\[
 p=\sum_{s\ {\rm small}}b_s,
 \qquad
 M_f=\prod_{s\ {\rm free}}(b_s+1).
\]
An empty product equals one.  If \(p=0\), we call
Proposition~\ref{prop:weight-oracle} once to estimate
\(h_H(\varnothing)=C(H)\), with relative tolerance \(\epsilon/32\)
and failure probability \(1/16\).  This branch, including the case of
an empty retained cell graph, is completed in the proof of
Theorem~\ref{thm:main} below.  Until then suppose that \(p\geq 1\).
All logarithms without a subscript are natural logarithms.

\subsection{Initialization and comparison of adjacent parameters}

\begin{lemma}[Exact initialization]
\label{lem:initialization}
For every state of the enlarged binary state space and
\(0\leq j<H\),
\begin{equation}
 \frac12\leq\frac{f_{j+1}}{f_j}\leq 1.
 \label{eq:alg-adjacent}
\end{equation}
Consequently the ideal multiplier \(C(j)/c_{il}(j)\) changes by a
factor in \([1/2,2]\) between consecutive parameters, and
\(\mu_j\leq 2\mu_{j+1}\) pointwise.  Moreover,
\begin{equation}
 \min_{S\in\mathcal A}\mu_j(S)
 \geq 2^{-(H+p+K)}
 \qquad (0\leq j\leq H).
 \label{eq:alg-minimum-mass}
\end{equation}
The values \(C(0)\) and \(c_{il}(0)\), and hence all initial ideal
multipliers, are computable exactly in polynomial bit time.  The law
\(\mu_0\) has the following finite-choice description: independently
for each small cell \(s\), choose a uniform integer
\(k_s\in\{0,\ldots,b_s\}\), and then a uniform \(k_s\)-element subset
of its \(b_s\) labeled pairs to use their row elements.
\end{lemma}

\begin{proof}
Lemma~\ref{lem:forest} gives \(0\leq D\leq H\).  Each summand in
\(h_{j+1}\) is therefore between one half and one times the
corresponding summand in \(h_j\), proving
\eqref{eq:alg-adjacent}.  The factorial factor in \(f_j\) is
independent of \(j\).  Summing the same inequalities over any type
gives
\[
 \frac12\leq \frac{C(j+1)}{C(j)}\leq 1,
 \qquad
 \frac12\leq \frac{c_{il}(j+1)}{c_{il}(j)}\leq 1.
\]
These bounds imply the multiplier comparison and
\(\mu_j\leq 2\mu_{j+1}\).

On a transversal the factorial factor is
\(\prod_s\binom{b_s}{x_s}^{-1}\geq 2^{-p}\).
Also \(h_j(x,y)\geq M_f2^{-H}\), while
\[
 C(j)\leq M_f\prod_{s\ {\rm small}}(b_s+1)
 \leq M_f2^K.
\]
Their quotient proves \eqref{eq:alg-minimum-mass}.

At parameter zero, \(h_0=M_f\), so
\begin{equation}
 C(0)=M_f\prod_{s\ {\rm small}}(b_s+1).
 \label{eq:alg-initial-total}
\end{equation}
We give an explicit formula for every defect total.  Fix distinct
pair labels \(i,l\), with \(i\) empty and \(l\) full.  For each small
cell \(s\), let
\[
 u_s=\mathbf 1_{\{i\text{ belongs to }s\}},\qquad
 v_s=\mathbf 1_{\{l\text{ belongs to }s\}},\qquad
 n_s=b_s-u_s-v_s.
\]
There are \(n_s\) singly occupied pairs left in that cell.  If exactly
\(k\) of these use their row elements, then
\(x_s=k+v_s\), \(y_s=n_s-k+v_s\), with multiplicity
\(\binom{n_s}{k}\).  Hence
\begin{equation}
 c_{il}(0)
 =M_f\prod_{s\ {\rm small}}
 \left[
 \frac{1}{b_s!}
 \sum_{k=0}^{n_s}
 \binom{n_s}{k}(k+v_s)!(n_s-k+v_s)!
 \right].
 \label{eq:alg-initial-defect}
\end{equation}
This includes the case in which \(i,l\) lie in the same cell:
then \(u_s=v_s=1\) and \(n_s=b_s-2\geq 0\).
All factorial arguments are at most \(p\), and all sums have at most
\(p+1\) terms.  Using the common denominator \(b_s!\) within each
bracket, the numerators and denominators in
\eqref{eq:alg-initial-defect} have
\(O(L+p\log(p+1))\) bits.  There are at most \(p(p-1)\) such
totals.  The multipliers
\[
 w_{il}^{(0)}=\frac{C(0)}{c_{il}(0)}
\]
are therefore positive rationals of polynomial bit length and are
computable by bounded integer arithmetic.
In fact the local bracket in \eqref{eq:alg-initial-defect} equals
\(b_s+1\), \(1\), \((b_s+1)(b_s+2)/6\), or
\((b_s+1)/6\), according as the cell contains neither designated
pair, only the hole, only the full pair, or both.  If \(s(i)\)
denotes the cell containing pair \(i\), the equivalent initializer is
\[
 w_{il}^{(0)}
 =
 \begin{cases}
 6,&s(i)=s(l),\\[2mm]
 \displaystyle\frac{6(b_{s(i)}+1)}{b_{s(l)}+2},
       &s(i)\ne s(l).
 \end{cases}
\]

Finally, the asserted sampler assigns any particular transversal
probability
\[
 \prod_{s\ {\rm small}}
 \frac{1}{(b_s+1)\binom{b_s}{x_s}}
 =\frac{f_0(S)}{C(0)}.
\]
This proves its distribution.  The finite-choice description will be
implemented with bounded fair-bit draws below.
\end{proof}

\subsection{A reference experiment with exact transition probabilities}

We first analyze a reference experiment using exact weights and
probabilities.  It is an auxiliary probability model; the implemented
algorithm in the next subsection uses only the explicit weight
evaluation procedure and independent fair bits.

To see which statistics a phase must collect, fix a parameter
\(0\le j<H\) and positive multipliers \(w\), and let \(\pi\)
be the stationary law of \(P_{j,w}\). Write \(p_T\) for the
stationary probability of type \(T\), where \(T=0\) denotes
\(\mathcal A\) and \(T=(i,l)\) denotes \(\mathcal A_{il}\).
Define
\[
 G_j(S)=\mathbf 1_{\mathcal A}(S)\frac{f_{j+1}(S)}{f_j(S)},
 \qquad u=\mathbb E_\pi G_j.
\]
The transversal mass is \(C(j)\), and the multiplied mass of
defect type \(il\) is \(w_{il}c_{il}(j)\). Therefore
\begin{equation}
 \frac{u}{p_0}=\frac{C(j+1)}{C(j)},\qquad
 w_{il}\frac{p_0}{p_{il}}=\frac{C(j)}{c_{il}(j)}.
 \label{eq:alg-observable-identities}
\end{equation}
Thus one trajectory can estimate both the next counting ratio and
the multipliers needed to balance the next phase. The type indicators
are constant on each defect class, and \(G_j\) vanishes there; all
these observables are fixed by \(\mathcal Q\), so the time-average
bound of Lemma~\ref{lem:time-average} applies when the multipliers
are good and the initial law is dominated by \(\pi\).

The stored multipliers depend on earlier trajectories. We therefore
start each phase with fresh randomness, using a new \(\mu_0\) draw
and trace stages at parameters \(1,\ldots,j\) with the multipliers
already learned. The actual restart has a deterministic transition
budget. For analysis we also consider the same restart without that
budget. Conditional on a past history in which the learned multipliers
are good, adjacent-parameter comparison and trace mixing give the
required domination for this virtual restart. Its cap failure is
bounded separately, without conditioning the endpoint law on meeting
the cap.

Let \(C_{\rm p},C_{\rm d},K_{\rm obs}\) have the values in
Section~\ref{sec:transport}, so in particular
\[
 K_{\rm obs}=10p^4(9C_{\rm p}+2C_{\rm d}).
\]
Choose the following integer time bounds and relative tolerance:
\begin{equation}
\begin{split}
 \eta&=\frac{\epsilon}{200(H+1)},\\
 \tau&=2K_{\rm obs}(H+p+K+5),\\
 B_{\rm ret}&=2000p^2(H+1)^2\tau,\\
 N_{\rm av}
 &=\left\lceil
 \frac{10^8p^8K_{\rm obs}(H+1)}{\eta^2}
 \right\rceil.
\end{split}
\label{eq:alg-phase-parameters}
\end{equation}
The quantities defining \(\tau\) and \(B_{\rm ret}\) are integers.
Initialize \(w^{(0)}\) by Lemma~\ref{lem:initialization}.  For
\(j=0,\ldots,H-1\), carry out the following phase.

\begin{enumerate}[label=\textbf{\arabic*.},leftmargin=*]
\item
\emph{Restart on transversals.}
Draw a fresh state from \(\mu_0\).  For \(a=1,\ldots,j\), run
\(\tau\) successive trace steps of \(P_{a,w^{(a)}}\), continuing
from the transversal reached at the preceding stage.  A trace step
ends at the next transversal after a strictly positive number of
underlying transitions.  Across this entire restart, allow at most
\(B_{\rm ret}\) underlying transitions.  If another transition is
needed after that budget is exhausted, abort the run and output zero.
For \(j=0\), the fresh \(\mu_0\) draw is the restart endpoint.

\item
\emph{Observe the current chain.}
From this endpoint, run \(P_{j,w^{(j)}}\) for \(N_{\rm av}\)
consecutive observed states, including the starting state.  For each
type \(T\), let \(n_T\) be its number of occurrences and
\(\widehat p_T=n_T/N_{\rm av}\).
Also form the empirical mean \(\widehat u\) of \(G_j\).
If any type count is zero, abort the run and output zero.

\item
\emph{Estimate the ratio and update the multipliers.}
Set
\begin{equation}
 \widehat R_j=\frac{\widehat u}{\widehat p_0}.
 \label{eq:alg-ratio-estimate}
\end{equation}
If \(j+1<H\), store
\begin{equation}
 w_{il}^{(j+1)}
 =w_{il}^{(j)}\frac{\widehat p_0}{\widehat p_{il}}
 =w_{il}^{(j)}\frac{n_0}{n_{il}}.
 \label{eq:alg-multiplier-update}
\end{equation}
\end{enumerate}

If no abort occurs, the reference output is
\[
 Z_{\rm ref}=C(0)\prod_{j=0}^{H-1}\widehat R_j.
\]
Every phase uses fresh randomness.  In particular, its initial
\(\mu_0\) sample and its chains do not reuse the final state of an
earlier observation trajectory.

\begin{proposition}[Accuracy of the reference experiment]
\label{prop:annealing}
With probability at least \(1-1/50\), the reference experiment does
not abort, every stored multiplier collection is good at its own
parameter in the sense of \eqref{eq:good-multipliers}, and
\begin{equation}
 \left|\log\frac{Z_{\rm ref}}{C(H)}\right|
 \leq 4H\eta.
 \label{eq:alg-reference-error}
\end{equation}
\end{proposition}

\begin{proof}
We make the conditioning in the argument explicit.  Let
\(\mathcal F_j\) be the sigma-field generated by the completed
trajectories, transition decisions, and counts of phases before
\(j\) in the capped reference experiment.  The stored multipliers
and the event that all these earlier phases were successful are
\(\mathcal F_j\)-measurable.  Fix a past history on that event.
The multipliers \(w^{(0)},\ldots,w^{(j)}\) are now fixed and good;
this will follow inductively from the update calculation below.
The current phase uses independent fresh randomness.

For analysis of this one phase, define a virtual experiment that
continues its restart without a transition cap, and then takes its
\(N_{\rm av}\) observations using fresh randomness.  Write
\(M=H+p+K\).  Lemma~\ref{lem:trace} and
\eqref{eq:alg-minimum-mass} imply that, from any starting
transversal, a stage of \(\tau\) trace steps has pointwise relative
error at most
\[
 2^M\exp(-\tau/K_{\rm obs})
 =2^M\exp(-2(M+5))<1.
\]
Its endpoint law is consequently at most \(2\mu_a\).
At the entrance to stage \(a=1\), the law is
\(\mu_0\leq 2\mu_1\).  At every later stage it is at most
\(2\mu_{a-1}\leq 4\mu_a\).  Invariance of the trace law preserves
the bound \(4\mu_a\) throughout that stage.

For good multipliers, \(\pi(\mathcal A)\geq 1/(5p^2)\).
The warm return-cost bound of Lemma~\ref{lem:trace} therefore
bounds the expected number of underlying transitions in each trace
step by \(20p^2\).  The virtual restart uses at most \(H\tau\)
trace steps, so its expected total cost is at most
\(20p^2H\tau\).  Markov's inequality gives
\begin{equation}
 \mathbb P\{\text{virtual restart cost}>B_{\rm ret}
                \mid\mathcal F_j\}
 \leq \frac{H}{100(H+1)^2}
 \leq \frac{1}{100(H+1)}.
 \label{eq:alg-cap-probability}
\end{equation}
For \(j=0\), the cost is zero and the same upper bound holds.

The virtual restart endpoint is at most \(2\mu_j\), including
\(j=0\).  Since \(\mu_j\) is the conditional law of the current
stationary distribution \(\pi\) on \(\mathcal A\), this endpoint
law is at most \(10p^2\pi\) on the full state space.
Each type indicator and \(G_j\) is fixed by the projection
\(\mathcal Q\) of Section~\ref{sec:transport}.  They take values
in \([0,1]\).  Their stationary means \(p_T,u\), given by
\eqref{eq:alg-observable-identities}, are at least \(1/(20p^2)\):
this is the type probability bound for defect types, while
\(p_0\geq 1/(5p^2)\) and \(C(j+1)/C(j)\geq 1/2\).

By Lemma~\ref{lem:time-average}, each empirical mean in the virtual
observation trajectory has mean squared error about its stationary
mean at most
\[
 \frac{20p^2K_{\rm obs}}{N_{\rm av}}.
\]
There are \(p(p-1)+2\leq 2p^2\) observables.  Markov's
inequality and a union bound therefore give
\begin{equation}
\begin{split}
 &\mathbb P\{\text{some virtual empirical mean has
 relative error greater than }\eta
               \mid\mathcal F_j\}\\
 &\hspace{12mm}\leq
 \frac{16000p^8K_{\rm obs}}{N_{\rm av}\eta^2}
 \leq\frac{1}{100(H+1)}.
\end{split}
\label{eq:alg-empirical-probability}
\end{equation}
These bounds concern the virtual, uncapped current phase.  Whenever
its restart cost is at most \(B_{\rm ret}\), it agrees with the
actual capped phase.  Thus failure of the capped phase is contained
in the union of the two virtual failure events in
\eqref{eq:alg-cap-probability} and
\eqref{eq:alg-empirical-probability}.  In particular, we never
condition the endpoint law on surviving the cap.

It remains to check that empirical success maintains the induction
and controls the output.  Because \(\eta<1/4\), empirical success
makes every type count positive. By
\eqref{eq:alg-observable-identities}, the exact ratio
\(w_{il}^{(j)}p_0/p_{il}\) is the current ideal multiplier
\(C(j)/c_{il}(j)\). The relative errors in
\(\widehat p_0,\widehat p_{il}\) change this value by a factor between
\((1-\eta)/(1+\eta)\) and its reciprocal, both lying in
\([1/2,2]\).  Lemma~\ref{lem:initialization} changes the ideal
multiplier by at most another factor two at the next parameter.
Equation~\eqref{eq:alg-multiplier-update} therefore satisfies
\eqref{eq:good-multipliers} for \(j+1\).  The update estimates the
current ideal multiplier directly; its error does not multiply the
earlier error in \(w^{(j)}\).

The same empirical bounds, and
\(u/p_0=C(j+1)/C(j)\), give
\[
 \left|\log\frac{\widehat R_j}{C(j+1)/C(j)}\right|
 \leq 4\eta.
\]
Initialization is good by Lemma~\ref{lem:initialization}.
For every successful fixed past, the conditional probability of a
current failure is at most \(2/[100(H+1)]\).  Summing the
probabilities of the first failing phase gives total failure
probability at most
\[
 \frac{2H}{100(H+1)}<\frac1{50}.
\]
On the complementary event, the ratio product telescopes and its
logarithmic errors sum to \eqref{eq:alg-reference-error}.
\end{proof}

\subsection{A bounded implementation using fair bits}

All parameters and counters below are represented exactly.  Set
\begin{equation}
 J_*=100(H+1)(B_{\rm ret}+N_{\rm av}+p+d+1),
 \qquad
 \xi=\frac{\eta}{1000J_*}.
 \label{eq:alg-bit-parameters}
\end{equation}
To evaluate any individual \(f_j(S)\), call
Proposition~\ref{prop:weight-oracle} with
\(\xi_{\rm in}=\theta=\xi\) and the small-count profile of \(S\),
using fresh bits on every call, and multiply by the exact factorial
factor in \eqref{eq:lifts}.  The evaluation is positive on every
execution.  For a valid proposed move, compute the ratio of the two
resulting rational \(\lambda\)-weights, clip it at one, and use a
dyadic acceptance threshold.  Use the analogous ratio of fresh
evaluations in each nonzero observation of \(G_j\).  Type counts and
the updates \eqref{eq:alg-multiplier-update} remain exact integer
and rational computations.  The restart transition caps are
unchanged.

For every finite-choice draw from \(\{0,\ldots,v-1\}\), apply
Lemma~\ref{lem:dyadic-integer-draw} with \(M=v\) and \(\zeta=\xi\).
This uses a fixed number of fair bits and has total variation error
at most \(\xi\).
For the subset draw in Lemma~\ref{lem:initialization}, take
\(v=\binom{b_s}{k_s}\) and decode an index in a fixed order of
the subsets.  Successive binomial inclusion/exclusion counts do
this with at most \(b_s\) decisions.  Thus the initial sampler uses
at most two finite-choice draws per small cell.  The choices of a
current element and a complementary element in a Metropolis proposal
likewise require only two scalar draws.

If \(\alpha\in[0,1]\) is a computed rational acceptance
probability, use
\(\lceil\log_2(1/\xi)\rceil\) fair bits and the downward-rounded
dyadic threshold for \(\alpha\).  Its probability differs from
\(\alpha\) by at most \(\xi\).  A single fair bit implements the
half hold exactly.  These procedures use fixed bit counts at their
requested resolutions and no rejection loop.

\begin{samepage}
\begin{lemma}[Coupling and bounded computation]
\label{lem:bounded-implementation}
The implemented run uses a number of bit operations bounded, on
every execution, by a fixed polynomial in \(L\) and
\(\epsilon^{-1}\).  It returns a nonnegative rational of polynomial
bit length.  It can be coupled to the capped reference experiment
so that, except on an event of probability at most \(10J_*\xi\),
all state trajectories, type counts, cap decisions, and stored
multipliers agree, every requested weight evaluation meets its
relative tolerance, and each positive observed term of \(G_j\)
differs from its reference value by a factor whose logarithm has
absolute value at most \(4\xi\).
\end{lemma}
\end{samepage}

\begin{proof}
We first bound the number of outer random actions.  Let
\(s_{\rm sm}\) be the number of positive small cells, so
\(s_{\rm sm}\leq\min(p,d)\), and put
\[
 T_{\max}=H(B_{\rm ret}+N_{\rm av}).
\]
This bounds all underlying transitions in the capped run.
Each transition needs at most two weight evaluations, four outer
draws (a hold coin, two element choices, and an acceptance threshold),
and one transition count.  Each of at most \(HN_{\rm av}\)
observations needs at most two additional weight evaluations.
Initial sampling needs at most \(2s_{\rm sm}\) outer draws per
phase.  Consequently
\begin{equation}
\begin{split}
 &\#\{\text{evaluation calls, outer draws, and transitions}\}\\
 &\quad\leq
 7T_{\max}+2HN_{\rm av}+2Hs_{\rm sm}\\
 &\quad\leq
 9H(B_{\rm ret}+N_{\rm av})+2H(p+d)<J_*.
\end{split}
\label{eq:alg-action-budget}
\end{equation}
A visited state has a single type, whose counter is incremented.
There is no separate sampling for each defect type.  Initializing
the \(p(p-1)\) counters and computing their updates are additional
deterministic polynomial computations.

For the coupling, expose the completed joint history of both
processes immediately before each requested evaluation or draw.
The event that the processes have matched so far, and the current
evaluation input, are measurable in this history.  A new evaluation
uses a fresh bit stream independent of that history.  Its conditional
failure probability is therefore at most \(\xi\), uniformly over
the requested profile.  This conclusion remains valid after earlier
matching and accuracy events, because they concern only already
exposed randomness.

Before a discrepancy, both runs have the same state and stored
multipliers.  Couple their initial and proposal draws at total
variation cost at most \(\xi\) per finite choice.  Conditional
on matching proposals and accurate evaluations, let \(r\) be the
reference acceptance ratio before clipping and \(r'\) the computed
ratio.  The exact factorial and multiplier factors are the same in
both runs, and hence
\[
 \frac{1-\xi}{1+\xi}\leq\frac{r'}r
 \leq\frac{1+\xi}{1-\xi}.
\]
It follows, regardless of the magnitude or quality of the
multipliers, that
\[
 \big|\min(1,r')-\min(1,r)\big|
 \leq \frac{2\xi}{1-\xi}<4\xi.
\]
The dyadic threshold adds at most \(\xi\) to this absolute
probability difference. To realize a common-uniform coupling, let
\(Y\) be the actual fresh uniform \(k\)-bit integer used for the
acceptance decision, and write its dyadic threshold as \(m/2^k\).
In the proof alone, introduce \(V\) uniform on \([0,1)\), independent
of \(Y\), the completed joint past, and the current weight evaluations.
Then \(U=(Y+V)/2^k\) is uniform on \([0,1)\), and the implemented
test \(Y<m\) is exactly \(U<m/2^k\). Use \(U<\min(1,r)\) for
the reference decision. Their mismatch probability is the absolute
difference of these thresholds, with the bound just obtained.
The auxiliary variable \(V\) is never sampled by the algorithm.

The coupling preserves the reference transition marginal
conditional on the completed joint past.  We do not condition the
reference process on the event that the entire coupling eventually
succeeds.  Matching paths have exactly the same integer type counts,
and the update \eqref{eq:alg-multiplier-update} depends only on
those counts.  Thus their stored multipliers remain exactly equal,
even though their observations of \(G_j\) need not be equal.
Trace return times and cap decisions are functions of the matched
state paths, so they also agree.  A union bound over
\eqref{eq:alg-action-budget}, charging evaluation failures and
draw or acceptance discrepancies, is at most \(10J_*\xi\).
On the remaining event, each positive observed weight ratio has
logarithmic relative error at most \(4\xi\).  Zero observations
are zero in both runs.

We next prove the unconditional resource bound.  Since all input
cells are listed explicitly and the parameters are included in
the input length, \(d,K=O(L)\).  The definitions in
\eqref{eq:parameters} give, for example,
\[
 A=O(L^3\epsilon^{-1}),\quad
 B=O(L^5\epsilon^{-1}),\quad
 W=O(L^6\epsilon^{-1}),\quad
 H=O(L^{11}\epsilon^{-2}),\quad
 p\leq dW.
\]
All quantities in \eqref{eq:alg-phase-parameters} and
\eqref{eq:alg-bit-parameters}, including \(\xi^{-1}\), are
therefore bounded by fixed polynomials in \(L,\epsilon^{-1}\).
Their rational encodings have polynomial length.  The state can be
stored using \(2p\) membership bits or \(p\) element labels;
the small counts and state type can be recovered by bounded scans.

There is also a uniform size bound for stored multipliers on
unsuccessful histories.  Let \(B_0=O(L+p\log(p+1))\) bound the
initial numerator and denominator lengths obtained from
\eqref{eq:alg-initial-defect}.  On every nonaborting history,
\[
 w_{il}^{(a)}
 =w_{il}^{(0)}
   \prod_{r=0}^{a-1}\frac{n_0^{(r)}}{n_{il}^{(r)}},
 \qquad
 1\leq n_0^{(r)},n_{il}^{(r)}\leq N_{\rm av}.
\]
Thus, even without fraction reduction, its numerator and denominator
each have at most
\begin{equation}
 B_0+H\left\lceil\log_2(N_{\rm av}+1)\right\rceil
 \label{eq:alg-multiplier-length}
\end{equation}
bits, up to a fixed additive constant.  A zero count causes an
abort before any such division.  This bound does not assume good
multipliers or accurate weight calls.

Proposition~\ref{prop:weight-oracle} supplies a deterministic
polynomial bound on the work and output length of each evaluation,
including its statistically unsuccessful executions.  More
explicitly, its ratio averages always lie in \([1/2,1]\) and its
rounded correction average always lies in \([1/16,8]\), so every
call at parameter \(j\) returns a number in
\[
 [\,M_f2^{-j}/16,\ 8M_f\,].
\]
Its bounded dyadic precision and bounded numbers of rational
operations give a polynomial representation length as well.
The inputs of this subroutine are \(j\), a valid small-count
profile, and the fixed tolerances; multiplier values are not
subroutine inputs.  Multiplication by the exact factorial factor
adds only \(O(p\log(p+1))\) bits.  Combining this bound with
\eqref{eq:alg-multiplier-length} proves that every Metropolis
ratio and threshold comparison has polynomial-sized operands.

For the output arithmetic, let \(B_f\) be an unconditional
polynomial bound on the lengths of evaluated \(f\)-weights.
Each observation ratio has \(O(B_f)\)-bit numerator and
denominator.  Adding \(N_{\rm av}\) such rationals using their
product as a common denominator requires
\(O(N_{\rm av}B_f+\log(N_{\rm av}+1))\) bits.  Dividing by
a nonzero empirical type probability adds only
\(O(\log(N_{\rm av}+1))\) bits.  The product of the \(H\)
phase estimates and \(C(0)\) therefore has polynomial bit
length.  Ordinary integer arithmetic on these bounded operands,
including the output operation, has polynomial cost.

Finally, every outer loop has its displayed deterministic cap,
and the evaluation subroutine itself has bounded loops and finite
bit draws.  Initialization and counter updates add at most
polynomially many operations on the bounded integers just described.
No part of this resource argument conditions on empirical accuracy,
on successful coupling, or on survival of a restart.  An abort
outputs the rational zero; otherwise all estimates are positive
rationals.  This proves the claim on every execution.
\end{proof}

\subsection{Relative error and amplification}

\begin{proof}[Proof of Theorem~\ref{thm:main}]
For an infeasible input the preliminary deterministic test returns
zero on every execution, as required.  Suppose that the input is
feasible.

First let \(p\geq 1\).  Combine
Proposition~\ref{prop:annealing} with
Lemma~\ref{lem:bounded-implementation}.  Except on an event of
probability at most \(1/50+10J_*\xi\), there is no abort,
reference and implemented paths agree, and every positive
observation differs from its reference counterpart by a factor
between \(\exp(-4\xi)\) and \(\exp(4\xi)\).  The same bounds
hold for sums of these positive observations.  Since type counts
agree exactly, each implemented ratio estimate differs from its
reference ratio by at most that same logarithmic amount.  The
implemented output \(\widehat Z\) consequently satisfies
\begin{equation}
 \left|\log\frac{\widehat Z}{C(H)}\right|
 \leq H(4\eta+4\xi)\leq\frac{\epsilon}{16}.
 \label{eq:alg-final-log-error}
\end{equation}
The failure bound is less than \(1/4\), since
\(10J_*\xi=\eta/100\).

Proposition~\ref{prop:contamination} gives
\[
 Z\leq C(H)\leq (1+\epsilon/32)Z.
\]
Together with \eqref{eq:alg-final-log-error}, this implies the
desired relative inequalities.  Indeed,
\(e^{-\epsilon/16}\geq 1-\epsilon/16\), and for
\(0<\epsilon<1\),
\[
 e^{\epsilon/16}(1+\epsilon/32)
 \leq (1+\epsilon/8)(1+\epsilon/32)
 \leq 1+\epsilon.
\]
Thus a single implemented run succeeds with probability greater
than \(3/4\).

The case \(p=1\) is already included in this argument.  The
enlarged state space then consists of the two transversals, the
defect-type family and multiplier updates are empty, and
\(\pi(\mathcal A)=1\).  Only the type-zero indicator and
\(G_j\) are observed.  All denominators and time bounds used
above remain defined.

If \(p=0\), there is a single small-count profile and
\(C(H)=h_H(\varnothing)\).  The call specified at the start of
this section returns a positive rational satisfying
\[
 (1-\epsilon/32)C(H)
 \leq \widehat Z
 \leq (1+\epsilon/32)C(H)
\]
with probability at least \(15/16\).  The contamination bound
again places this interval inside
\([(1-\epsilon)Z,(1+\epsilon)Z]\).  The evaluation procedure
covers zero free coordinates as well, in which case its sum and
bin space each have one element.  This branch invokes none of
the \(p\)-dependent chain parameters and has the same required
polynomial bit bound.

It remains to obtain the requested failure probability.  Put
\[
 s=\left\lceil\log_2(\delta^{-1})\right\rceil,\qquad
 R=10s+1,
\]
run the appropriate constant-success algorithm independently
\(R\) times, and return the median of its nonnegative rational
outputs.  Let \(F\) be the number of runs outside the desired
relative interval.  Independence and the single-run failure
bound \(1/4\) give
\[
 \mathbb E\,2^F\leq(5/4)^R.
\]
A bad median requires at least half the runs to fail, so
\[
 \mathbb P\{\text{bad median}\}
 \leq \left(\frac{5}{4\sqrt2}\right)^{10s+1}
 \leq 2^{-s}\leq\delta.
\]
The middle inequality follows, for example, from
\((5/(4\sqrt2))^{10}<1/2\).

Lemma~\ref{lem:bounded-implementation}, or
Proposition~\ref{prop:weight-oracle} for \(p=0\), bounds the work
and output length of every run by one fixed polynomial in
\(L,\epsilon^{-1}\).  The number of repetitions is
\(O(\log(\delta^{-1}))\); exact comparisons and median selection
on their polynomial-length rational outputs also have polynomial
bit cost.  Ceilings of binary logarithms used for the parameters
and bit counts are computable by integer shifts and comparisons.
All randomness in the implemented procedure consists of
independent unbiased bits.  Partition functions, defect totals
away from parameter zero, proof-only conditioning trees, and
transport flows are never queried or enumerated.  The only weight
evaluation used is the bounded procedure proved in
Proposition~\ref{prop:weight-oracle}.  This establishes a uniform
FPRAS with the stated bound on every execution.
\end{proof}
